\documentclass[11pt,letterpaper]{article}
\usepackage{technical-design-spec}
\usepackage{longtable,tabularx,array}
\setDocTitle{Data Governance Standards and Frameworks Technical Reference}
\setDocSubtitle{Enterprise reference and control architecture}
\setDocVersion{1.0}
\setDocStatus{Proposed / original reference}
\setDocOwner{Enterprise Data Governance}
\setDocDate{30 October 2026}
\setDocSummary{Original implementation guidance; source standards and official assessments remain authoritative.}
\begin{document}
\makedoctitle
\makedocsummary
\noindent\textbf{Authority boundary.} Independently authored guidance; not a reproduction of licensed requirements, legal advice or an official DCAM assessment.
\section{Authority taxonomy}
External data governance references have different authority: \textbf{international standards} define particular requirements or guidance; \textbf{professional bodies of knowledge} provide practices and vocabulary; \textbf{governance frameworks} structure decision processes; \textbf{capability models} assess adoption; \textbf{laws} impose entity- and activity-dependent obligations. Internal enterprise policy and contracts may independently make practices mandatory. Never equate adoption with certification, or a publisher overview with a clause-level compliance assessment.

\section{Authoritative reference landscape}
\begin{longtable}{p{.25\linewidth}p{.34\linewidth}p{.33\linewidth}}
\textbf{Reference} & \textbf{Primary scope} & \textbf{Implementation artifact}\\\hline
DAMA-DMBOK & Enterprise data management across eleven knowledge areas & Data strategy, domains, steward charter, glossary\\
ISO/IEC 38505-1:2026 & Governing-body governance of organizational data & Oversight charter, delegated decisions and review minutes\\
ISO 8000-1:2022 & Data quality series overview & Approved quality policy and measurement taxonomy\\
ISO 8000-51:2023 & Exchange of governance policy statements & Machine-readable policy statements and checks\\
ISO/IEC 11179-3:2023 & Metadata registry common facilities & Versioned terms, definitions, mappings and registration\\
ISO/IEC 11179-31:2023 & Data specification registration & Data element, domain and permissible value register\\
ISO/IEC 5259-2:2024 & Analytics and ML data quality measures & Dataset fitness metrics and thresholds\\
ISO/IEC 5259-5:2025 & Oversight of analytics and ML quality & Lifecycle governance review and owners\\
COBIT 2019 & Enterprise governance/management of information and technology & Approved controls, role assignment, assurance\\
EDM DCAM v3 & Evidence-based data management capability assessment & Engagement, process and evidence records\\
\end{longtable}
\textbf{Interpretation caveat:} this is a local engineering crosswalk and does not assert equivalence between requirements. DCAM v3 identifies eight components, 34 capabilities and 101 sub-capabilities; official instrument content and scoring are not reproduced.

\section{Governance operating layers}
A practical control system separates governing body direction, management accountability, semantic ownership, technical implementation, evidence collection and independent assurance.
\begin{enumerate}
\item Establish business outcomes, risk appetite, jurisdiction and system boundaries.
\item Classify the source by authority; record publisher, edition, part, licence, effective date and reviewer.
\item Approve applicability and a plain-language enterprise control objective.
\item Assign exactly one accountable owner, responsible operator, test owner and risk exception authority.
\item Map policy and business glossary IDs to conceptual model, logical attributes, physical schema and lineage.
\item Define test method, population, sampling frequency, threshold, retention and evidence location.
\item Test operating effectiveness and record deficiency, exception, remediation and independent sign-off.
\item Re-evaluate after material regulatory, standards, processing or system changes.
\end{enumerate}

\section{Selection and traceability method}
For non-mandatory reference selection, evaluate business coverage (25 percent), authority fit (25 percent), implementability (20 percent), evidence measurability (15 percent), and maintenance (15 percent). Rate each from 0 to 4, calculate \(100\sum w_i r_i/4\) with weights as decimal fractions, and record alternatives and the final decision. These are locally proposed weights, \emph{not} ISO or DCAM scoring. A mandatory legal obligation is not waived by a low suitability score.
Maintain stable chains: external source ID and edition \(\rightarrow\) interpreted requirement \(\rightarrow\) enterprise policy \(\rightarrow\) testable control \(\rightarrow\) implementation \(\rightarrow\) secured evidence \(\rightarrow\) assessment finding and closure. Missing evidence is not an automatic pass.

\section{Data quality, metadata and architecture}
For a critical data element, register an approved business definition, domain and conceptual entity, logical attribute, physical column or event field, valid value set, owner and steward, source-to-target lineage, privacy class and retention disposition. A quality rule should specify dimension, numerator and denominator, sample window, measurable threshold and remediation authority. Connect schema and data contracts to stewardship and controlled metadata IDs. Verify actual production behavior rather than accepting a design diagram as operating evidence. See \texttt{src/data-systems/enterprise-data-modeling/}.

\section{Regulatory applicability}
GDPR, the EU Data Governance Act, California CCPA/CPRA, HIPAA, GLBA, SOX and UK data protection legislation do not apply uniformly to every enterprise. Record territorial nexus, organizational role, legal basis or activity, data categories, obligations, exclusions, authoritative citation, amendment, effective date and legal approval. The EU Data Governance Act applies from 24 September 2023 and does not replace GDPR. The UK Data (Use and Access) Act 2025 amended existing legislation, with affected data protection provisions in force by June 2026. Legal review is required before asserting applicability.

\section{Metrics, evidence and review}
Report critical-data ownership coverage, approved definition coverage, data quality pass rates, lineage coverage, access review completion and tested control coverage. Publish numerator, denominator, population, date, method, unresolved exceptions and confidence limitations. Evidence must include provenance, controlled access, integrity where needed, retention, test procedure, result, reviewer and change history. Audit findings have an accountable owner and dated retest; a link alone proves nothing.

\section{Publisher source register}
\begin{itemize}
\item ISO/IEC 38505-1:2026: \url{https://www.iso.org/standard/87195.html}
\item ISO 8000-1:2022: \url{https://www.iso.org/standard/81745.html}
\item ISO 8000-51:2023: \url{https://www.iso.org/standard/78708.html}
\item ISO/IEC 11179-3:2023: \url{https://www.iso.org/standard/78915.html}
\item ISO/IEC 11179-31:2023: \url{https://www.iso.org/standard/78925.html}
\item ISO/IEC 5259-2:2024: \url{https://www.iso.org/standard/81860.html}
\item ISO/IEC 5259-5:2025: \url{https://www.iso.org/standard/5259-5}
\item ISACA COBIT: \url{https://www.isaca.org/resources/cobit}
\item EDM Association DCAM: \url{https://edmcouncil.org/frameworks/dcam/}
\item DAMA: \url{https://www.dama.org/}
\item EU DGA summary: \url{https://eur-lex.europa.eu/legal-content/EN/LSU/?uri=CELEX:32022R0868}
\item UK Information Commissioner's Office: \url{https://ico.org.uk/about-the-ico/what-we-do/legislation-we-cover/data-use-and-access-act-2025/}
\end{itemize}
Sources checked 30 October 2026. Normative clause-by-clause use requires the authorized official publications.

\end{document}